% ECONOMICS_PROBLEM: {"id": "Q52-398", "title": "Distributional incidence of climate policy", "jel_code": "Q52", "source_id": "OP-0119"}
% !TeX program = xelatex
\documentclass[11pt,a4paper]{article}
\usepackage[margin=23mm]{geometry}
\usepackage{fontspec}
\setmainfont{Latin Modern Roman}
\usepackage{amsmath,amssymb,mathtools}
\usepackage{xcolor,enumitem,booktabs,longtable,array,xurl,hyperref}
\definecolor{muted}{HTML}{53636C}
\hypersetup{colorlinks=true,urlcolor=blue,linkcolor=blue}
\setlength{\parindent}{0pt}
\setlength{\parskip}{5pt}
\newcommand{\R}{\mathbb R}
\newcommand{\E}{\mathbb E}
\newcommand{\N}{\mathbb N}
\newcommand{\ind}{\mathbf 1}
\newcommand{\Var}{\operatorname{Var}}
\newcommand{\Cov}{\operatorname{Cov}}
\newcommand{\supp}{\operatorname{supp}}
\DeclareMathOperator*{\argmin}{arg\,min}
\DeclareMathOperator*{\argmax}{arg\,max}
\newcommand{\entrymeta}[3]{{\sffamily\small\color{muted}\textbf{\texttt{#1}}\quad|\quad #2\par #3\par}}
\newcommand{\Problemref}[1]{\texttt{#1}}
\newcommand{\JELref}[2]{\texttt{#2} (JEL \texttt{#1})}
\begin{document}
\section*{Distributional incidence of climate policy}
\textbf{Problem ID:} \texttt{Q52-398}\quad \textbf{JEL:} \texttt{Q52}\par
% BEGIN ECONOMICS STATEMENT
\label{problem:OP-0119}
\label{atlas:prob:C9}

\noindent\textbf{Original identifier:} C9 (atlas item 119). \textbf{Field:} Environmental and climate economics. \textbf{Classification:} editorial research family with established restricted solutions; current open status not certified.

\subsection*{Definitions and assumptions}

Let climate-policy regime $r\in\{0,1\}$ specify carbon prices, subsidies, financing and transfers over dates $t=0,\ldots,H$. A model $m$ and scenario $z$ yield equilibria $e_r\in\mathcal E(m,z,r)$ with goods prices, wages, rents, household asset values and emissions. Household $h$ has intertemporal preferences $U_h$, initial endowments, employer links and portfolio positions. Fix a real date-zero consumption numeraire. Define compensating variation $CV_h(e_0,e_1)$ as the date-zero transfer satisfying $U_h(e_1;CV_h)=U_h(e_0;0)$, if a unique solution exists, with equilibrium prices in $e_1$ held fixed for this money-metric comparison.

\subsection*{Formal research question}

Characterize identification and distributional inference for full policy incidence after product, labor, housing and asset-market adjustments. For a specified household group $G$ and weights $w_h\ge0$ summing to one within it, define burden $B_G=\sum_{h\in G}w_hCV_h$. Over models $\mathcal M(P)$ consistent with linked consumption, employment and balance-sheet data $P$, estimate or bound the joint vector of group burdens and its distribution. State how equilibria are selected; absent selection, project all admissible equilibrium pairs into a burden set. Determine which policy shocks separately identify consumption-price, factor-income, capitalization and revenue-recycling channels, and derive confidence regions separating sampling uncertainty from preference, model and scenario ambiguity.

\subsection*{Interpretation and scope}

Positive $CV_h$ denotes the compensation needed for a loss, while negative values denote gains. A static expenditure share captures only one channel. Employers may adjust wages or employment, housing markets capitalize location changes, and asset owners can gain or lose wealth. Capitalized price changes and future income changes must not be added twice; policy financing and government liabilities must satisfy the model's budget constraints. Ex ante welfare under risk requires a stated probability law and discounting, rather than evaluating only mean income. Incidence of a small transfer at fixed equilibrium prices differs from re-solving the economy after a large aggregate compensation program. Household grouping and population weights affect reported inequality.

\subsection*{Source audit and status}

[S1] studies the direct distribution of clean-energy tax credits, explicitly limiting broader incidence. The undated Cronin item is completed with the verified 2019 tax-and-rebate article [S2]. The undated Metcalf--Stock item is completed with the verified 2023 European-carbon-tax article [S3], distinct from their 2020 article; intended identities are uncertain. Aggregate macroeconomic effects do not settle household incidence. The full-equilibrium distributional target is editorial and these sources do not certify its exact current open status.

\subsection*{References}

\begin{enumerate}[label={[\textnormal{S}\arabic*]},leftmargin=*]

\item Severin Borenstein and Lucas W. Davis (2016). \emph{The Distributional Effects of US Clean Energy Tax Credits}. Tax Policy and the Economy 30(1), 191–234. \url{https://doi.org/10.1086/685597}. \textbf{Locator:} Publisher or institutional abstract and bibliographic record. \textbf{Support and limits:} Direct distribution of tax-credit receipt; expressly does not measure the full source-of-funds or general-equilibrium incidence.

\item Julie Anne Cronin, Don Fullerton and Steven Sexton (2019). \emph{Vertical and Horizontal Redistributions from a Carbon Tax and Rebate}. Journal of the Association of Environmental and Resource Economists 6(S1), S169–S208. \url{https://doi.org/10.1086/701191}. \textbf{Locator:} Publisher or institutional abstract and bibliographic record. \textbf{Support and limits:} Relevant verified completion of incomplete input. Household expenditure/recycling simulation does not identify all asset and factor-price channels.

\item Gilbert E. Metcalf and James H. Stock (2023). \emph{The Macroeconomic Impact of Europe’s Carbon Taxes}. American Economic Journal: Macroeconomics 15(3), 265–286. \url{https://doi.org/10.1257/mac.20210052}. \textbf{Locator:} Publisher or institutional abstract and bibliographic record. \textbf{Support and limits:} Selected verified candidate, distinct from their 2020 AEA Papers and Proceedings article. Aggregate GDP/employment findings do not settle distributional incidence.

\end{enumerate}

\subsection*{Common conventions: Source mathematical conventions}

Unless an entry states otherwise, agent and firm sets are finite. State and action spaces are measurable, strategies and outcome maps are measurable, probability laws are specified, and expectations exist for the targets under discussion. Infinite-horizon discount factors lie in $(0,1)$ and discounted objectives are finite. Norms, tolerances and complexity input sizes must be specified for computational comparisons. Constants or primitives that appear locally are inputs, not empirical facts asserted by this catalogue.

\textbf{Identification.} A statistical model $m\in\mathcal M$ implies an observable law $P_m$ and target $\theta(m)$. At law $P$, its identified set is
\[
\Theta(P;\mathcal M)=\{\theta(m):m\in\mathcal M,\ P_m=P\}.
\]
Point identification holds when this set is a singleton, or equivalently when $P_{m_1}=P_{m_2}$ implies $\theta(m_1)=\theta(m_2)$ for all compatible models. Sharp bounds mean bounds attained, or approached when the boundary is not attained, by compatible models. Writing this set is a definition, not a solution: the substantive task is to establish informative restrictions and derive or estimate the set. An unrestricted-confounding model may give only trivial bounds.

\textbf{Inference.} Identification, estimability and valid uncertainty quantification are separate questions. Fix $\alpha\in(0,1)$. A confidence procedure $C_n$ over a declared law class $\mathcal P$ has asymptotic uniform coverage at level $1-\alpha$ when
\[
\liminf_{n\to\infty}\inf_{P\in\mathcal P}\Pr_P\{\theta(P)\in C_n\}\geq1-\alpha.
\]
For set-identified targets the coverage event must be defined separately, for example coverage of the full identified set. If an entry uses identification-indexed models instead of uniquely indexed laws, the same coverage convention is applied over those models. Pointwise limits do not establish uniform validity, and asymptotic validity does not guarantee finite-sample precision. Sample sizes, dependence structures and weak-identification sequences are part of each application.

\textbf{Counterfactuals.} $Y_i(a)$ is a potential outcome under a specified intervention, including its timing, implementation and versions. It is not $Y_i$ conditional on voluntarily choosing $a$. A full intervention vector permits interference; shorthand scalar treatment notation does not silently impose no interference. Assignment, selection, attrition, anticipation and measurement must be specified. A claim about an instrument's local effect is not automatically a population policy effect.

\textbf{Economic equilibrium.} A counterfactual model includes best responses, constraints, feasibility, market clearing, state transitions and expectations consistent with the stated information. When several equilibria exist, either a declared selector is an input or the target is set-valued. Comparisons across policy regimes must use the stated selection convention. Reduced-form relationships are not presumed invariant after an equilibrium-changing intervention.

\textbf{Optimization and welfare.} Optimization is over the stated feasible set. If attainment is not established, $\sup$ or $\inf$ and $\varepsilon$-optimal choices replace an asserted maximizer or minimizer. For example, with value $V=\sup_{a\in\mathcal A}W(a)<\infty$, an $\varepsilon$-optimal set is $\{a\in\mathcal A:W(a)\geq V-\varepsilon\}$ for $\varepsilon>0$. Benefits, costs, risk and health or environmental outcomes can be aggregated only using declared units and conversion weights. Interpersonal utility weights, the treatment of fiscal transfers and foreign profits, discounting, and the behavioral welfare criterion are explicit inputs. A comparative-static derivative additionally requires existence and appropriate differentiability of the selected counterfactual.

\textbf{Sources.} Local labels [S1], [S2], and so forth restart in every question. Locators identify the relevant abstract, model, section or result at the level actually checked; a bibliographic or abstract-level verification is not represented as inspection of an entire proof. Working papers are distinguished from later publications and changing versions. Source summaries are paraphrased. All URL checks and editorial formulations refer to the audit date stated above.
% END ECONOMICS STATEMENT
\end{document}
